\begin{abstract}
Network sockets have been the dominant way for applications to interact with the
network for several decades. The design of the BSD socket API forces applications to
bind to an underlying transport protocol at compile time, which, together
with middlebox interference greatly hinder innovation at the transport layer.

The Transport Services (TAPS) working group has produced a series of
RFCs~9621-9623, describing an abstract, protocol-agnostic and asynchronous networking library. 
While previous implementations of TAPS exist, none have been developed since the final
release of RFC~9621, RFC~9622 and RFC~9623 in 2025. Additionally, no current implementations
of TAPS support QUIC, making it unclear exactly how QUIC, as the most successful recent protocol innovation
maps onto the TAPS API.

This thesis presents CTaps, a C implementation of IETF Transport Services with support for QUIC, TCP and
UDP. CTaps supports several important QUIC features, such as multistreaming, connection
migration and session resumption, all through the TAPS API as described in
RFC~9622. In addition, CTaps supports important TAPS features such as
candidate gathering, and candidate racing.
We find that QUIC maps well, with most
features either being directly controllable or inferable through the TAPS API.
Remaining feature gaps are minor, such as unidirectional stream resets.
We evaluate CTaps against native TCP and Picoquic QUIC clients across handshake
latency, multistream transfer, candidate racing and connection migration.
We argue that CTaps effectively serves as a reference for future implementations of TAPS,
especially when supporting QUIC.
\end{abstract}
